\section{Introduction}
\label{sec:intro}

Deploying large language models under memory constraints requires loading model parameters
from host memory at each forward pass.
For a 7B-parameter model in fp16, streaming roughly 14\,GB of weights from host memory at
16\,GB/s would impose an idealized lower bound of about 0.9\,s/token before compute and
overheads --- a regime in which memory movement, not arithmetic, is the binding constraint.
Quantization compresses individual parameters but preserves the all-weights-per-token access
pattern.
Mixture-of-Experts routing \citep{shazeer2017moe,fedus2022switch} reduces \emph{which}
parameters are active per token, but standard MoE architectures activate different expert
sets per layer, accumulating a transfer budget proportional to depth.

We study a complementary structure: a single shared \emph{block bank} that a recursive
router accesses repeatedly.
At each recursion step $r$, the router selects a small set of $k$ blocks from a bank of $n$
blocks and applies them to the running hidden state.
The key constraint --- the \textbf{SR-Core} invariant --- is that routing happens only at
$r = 1$; steps $r = 2, \ldots, R$ reuse the same block selection.
The active weight set per token is therefore exactly $k$ blocks, regardless of recursion
depth $R$.
We refer to this as the \textbf{working set} ($\WS = k$), and it is independent of bank
size $n$ and depth $R$.

This guarantee is architecturally meaningful: at inference time the identity of the $k$
active blocks is known before any matrix multiply is executed, enabling block-granular
prefetching.
With $n = 64$ blocks and $k = 8$, the working set covers $12.5\%$ of the bank.
Offloading simulation on trained models confirms the structural advantage: sparse SR-Core
loads approximately $4.1\times$ fewer bytes per token than dense layer-offloading at a
realistic LRU cache capacity of $K = 8$ blocks ($3{,}035$ vs.\ $12{,}348\kb$ at 16\,GB/s).
This is a simulated bytes-in-motion estimate under an LRU caching model, not a measured
latency speedup.

$\WS = k$ is, however, a \emph{necessary} condition for cache efficiency, not a sufficient
one.
Whether the actual bytes-in-motion are low depends on how the router distributes selections
across tokens.
If every token activates a different set of $k$ blocks, the effective working set
\emph{across a sequence} approaches $n$, and a small cache will miss on nearly every
request.
Cache efficiency requires that routing patterns are \textbf{locally stable} --- that nearby
or semantically similar tokens tend to route through overlapping block sets, so that a
modest cache serves a disproportionately large fraction of requests.

We show that unconstrained SR-Core training produces a router that is only partially stable,
and that a targeted entropy-based objective can systematically move the router along a
\textbf{cache/locality axis} --- trading a small amount of per-token routing diversity for
significantly improved block-level cache hit rates.
Mild entropy pressure improves held-out code behavior with modest cache gains; stronger
pressure consistently improves cache metrics, with quality effects that are seed-dependent.
Dense models remain the quality upper bound throughout; the contribution is a controllable
systems-level degree of freedom within the sparse model family.

\paragraph{Contributions.} We present:
\begin{enumerate}
  \item An analysis of earlier recursive-routing dynamics that motivated the SR-Core
        constraint, revealing a stable two-phase structure (token-specific entry; collapsed
        reusable core) and an inherent reuse/novelty tension.
  \item An entropy-based router consolidation objective that converts this tension into a
        tunable cache/locality axis, with methodologically consistent negative controls.
  \item Empirical evaluation on a four-domain pretraining benchmark confirming a
        reproducible Pareto frontier between cache efficiency and held-out code
        generalization, replicated across seeds, and bounded above by a dense quality
        baseline.
\end{enumerate}
